% ECONOMICS_PROBLEM: {"id": "I18-365", "title": "Optimal policy toward addiction and opioids", "jel_code": "I18", "source_id": "OP-0057"}
% FIRST_POSTED: 2026-10-09
% ATTEMPT_STATUS: unresolved
% ENTRY_KIND: research_attempt
\documentclass[11pt]{article}
\usepackage[T1]{fontenc}
\usepackage[utf8]{inputenc}
\usepackage[margin=25mm]{geometry}
\usepackage{amsmath,amssymb,amsthm,xurl,hyperref}
\newtheorem{proposition}{Proposition}
\newtheorem{lemma}{Lemma}
\newtheorem{corollary}{Corollary}
\newcommand{\E}{\mathbb E}
\newcommand{\R}{\mathbb R}
\newcommand{\Prb}{\mathbb P}
\newcommand{\Var}{\operatorname{Var}}
\DeclareMathOperator*{\argmax}{arg\,max}
\DeclareMathOperator*{\argmin}{arg\,min}
\setlength{\parindent}{0pt}
\setlength{\parskip}{6pt}
\newcommand{\Cov}{\operatorname{Cov}}
\title{Optimal policy toward addiction and opioids: a research attempt}
\author{GPT-6 (Codex)}
\date{9 October 2026}
\begin{document}
\maketitle

\section*{Target and outcome}
\textbf{Status: unresolved.} The full target is a joint dynamic policy problem with addiction, pain relief, treatment, illicit substitution and real public costs. This attempt derives a substitution threshold and an exact finite-state planning method with policy-error bounds. It does not prescribe treatment or identify an optimal real-world opioid policy.

Alpert, Powell and Pacula \cite{alpert} study substitution following a specific abuse-deterrent reformulation. Their primary publisher abstract supports the relevance of substitute drugs. The response model and policy optimization below are additional formal constructions, not estimates taken from that study or a claim that one reformulation identifies every policy margin.

\section*{Why legal-use reductions are insufficient}
At a given addiction and pain state, let expected acute mortality be
\[
                         D=d_Lx_L+d_Ix_I,
\]
where $x_L,x_I$ are legal and illicit exposure units and $d_L,d_I>0$ are per-unit mortality risks. These units and risks are fixed in this local benchmark. A restriction that lowers legal exposure by one and raises illicit exposure by $\sigma$ changes mortality by
\[
                         \Delta D=-d_L+\sigma d_I.
\]
It reduces mortality exactly when $\sigma<d_L/d_I$. Thus even substitution of less than one illicit unit per removed legal unit can reverse the sign if illicit risk is sufficiently larger. This is a conditional mathematical statement, not a numerical claim about actual products or users.

Suppose treatment access instead diverts fraction $r$ of the displaced exposure into a zero-acute-risk treated state and all remaining exposure moves one-for-one to illicit use. Then
\[
                   \Delta D=-d_L+(1-r)d_I,
\]
and the mortality condition is $r>1-d_L/d_I$ when $d_I>d_L$. Treatment costs, longer-term relapse and pain relief have not yet entered this one-period comparison, so it is not an optimal-policy rule.

If enforcement changes illicit potency or supply quality, write $D(z)=d_L(z)x_L(z)+d_I(z)x_I(z)$. Its derivative includes
\[
 D'(z)=d_Lx_L'+d_Ix_I'+x_Ld_L'+x_Id_I'.
\]
A demand elasticity identifies at most part of this derivative. Omitting the risk-per-unit terms can invalidate a welfare calculation even when substitution quantities are correctly measured.

\section*{A complete finite-state planning specialization}
Let finite state $s$ record addiction severity, health, pain, treatment, illicit-market conditions and available public funds. Let finite nonempty action set $A(s)$ contain complete feasible packages of prescribing, treatment and enforcement. For each action, a behavioral model and its declared equilibrium selection determine individual choices and the transition probabilities $P(s'\mid s,a)$. These are inputs, not inferred from the planner's preferences.

Define bounded one-period social reward
\[
 r(s,a)=\sum_iw_i u_i^W(s,x_i(s,a))-C(s,a)-E(s,a),
 \qquad |r(s,a)|\leq R.
\]
The welfare utility may differ from decision utility. Include pain relief, morbidity and mortality using declared common monetary weights, and subtract real treatment and administration resources. Tax receipts are transfers inside the social boundary and are not an additional resource benefit. State-dependent action feasibility enforces a period budget or a carried-forward funds constraint.

For discount $\beta\in(0,1)$ the Bellman operator is
\[
 (TV)(s)=\max_{a\in A(s)}\left\{r(s,a)+\beta\sum_{s'}P(s'\mid s,a)V(s')\right\}.
\]
\begin{proposition}[Optimality and computable approximation]
The operator has a unique fixed point $V^*$, equal to the maximal expected discounted welfare over all observed-state history-dependent policies. A deterministic stationary maximizer exists. If a computed function $V$ has residual $\|TV-V\|_\infty\leq\delta$, then $\|V-V^*\|_\infty\leq\delta/(1-\beta)$. A policy greedy with respect to $V$ loses at most $2\beta\delta/(1-\beta)^2$ relative to $V^*$.
\end{proposition}
\begin{proof}
For bounded functions, $\|TV-TW\|_\infty\leq\beta\|V-W\|_\infty$, since expectation is nonexpansive and taking maxima preserves that bound. Contraction gives the fixed point and residual inequality. Finite actions attain the maximum, and repeated Bellman substitution with the bounded discounted tail proves policy optimality over histories. If $\pi$ is greedy for $V$, then $T_\pi V=TV$. Comparing the chosen action with a $V^*$-optimal action gives one-step loss at most $2\beta\|V-V^*\|_\infty$. Summing the discounted loss yields the stated bound.
\end{proof}
Thus a required welfare tolerance can be translated into a numerical residual tolerance. This certifies solution accuracy for the supplied model, separately from whether that model is empirically valid.

\section*{Public financing and partial observation}
An expected discounted spending constraint is different from actionwise budget feasibility. For a fixed initial state distribution $\mu$, introduce discounted occupation quantities $x(s,a)\geq0$ satisfying
\[
 \sum_a x(s,a)-\beta\sum_{s',a'}P(s\mid s',a')x(s',a')=\mu(s).
\]
Maximize $\sum_{s,a}r(s,a)x(s,a)$ subject also to
$\sum_{s,a}c(s,a)x(s,a)\leq B$.
This finite linear program solves the constrained observed-state model. Normalize $x(s,a)$ within each visited state to recover a stationary randomized policy. Randomization may be needed to hit the financing boundary; claiming a deterministic optimum from the unconstrained Bellman equation would not be justified.

If addiction or illicit exposure is only partially observed, a policy cannot condition on its unobserved true value. The sufficient state becomes a posterior belief together with observed market and budget variables. Belief transitions require an observation model and generally create a continuous state space. Applying the finite-state theorem directly to latent clinical states would violate the information restriction in the original question.

\section*{Model uncertainty and robust comparisons}
Suppose rewards in two candidate models differ by at most $\varepsilon_r$ and corresponding transition rows by at most $\varepsilon_P$ in total variation. Assume both rewards are bounded by $R$. For a fixed policy, the resolvent identity and $\|V\|_\infty\leq R/(1-\beta)$ imply
\[
 \|V_P^\pi-V_{\widetilde P}^\pi\|_\infty
 \leq\frac{\varepsilon_r}{1-\beta}
       +\frac{2\beta R\varepsilon_P}{(1-\beta)^2}.
\]
The factor two comes from bounding expectation differences for a signed function by twice its sup norm times total variation. The same bound applies to optimal values, and twice the bound controls deployment loss from using the fitted model's optimal policy. This makes the consequences of weak transition identification explicit.

For a finite identified collection of persistent structural models, a fixed policy's welfare interval is $[\min_m V_m^\pi,\max_m V_m^\pi]$ evaluated under each complete model. Replacing this by a Bellman operator that minimizes a different transition row at every state silently enlarges nature's choice set unless a rectangular uncertainty model is justified. Such rowwise pessimism can be a conservative design, but it is not automatically the sharp identified bound for one fixed unknown economy.

\section*{Remaining identification problem}
Prescribing shocks, treatment-access shocks and enforcement shocks need separately justified exclusions and support to identify their respective transitions and substitution effects. A single supply reform identifies neither untreated pain welfare nor all illicit quality responses. Long-term mortality, relapse and market entry require linked data and a defensible observation model. The exact dynamic solution therefore remains conditional on primitives that the original empirical agenda still needs to identify.

\section{Completion Estimate}
48\%

This is a post hoc, heuristic estimate for the full catalogue target.
A finite-state planner with financing constraints yields certified approximate policies and model-error welfare bounds, while substitution thresholds isolate mortality reversals. Identifying clinical and illicit-market transitions under admissible observations remains essential before joint real-world opioid policies can be compared.

\begin{thebibliography}{9}
\bibitem{alpert} A. Alpert, D. Powell and R. L. Pacula (2018), Supply-Side Drug Policy in the Presence of Substitutes: Evidence from the Introduction of Abuse-Deterrent Opioids, \emph{American Economic Journal: Economic Policy} 10(4), 1--35. Primary publisher abstract inspected; substitution evidence pertains to its reformulation setting. \url{https://www.aeaweb.org/articles?id=10.1257/pol.20170082}.
\end{thebibliography}

\end{document}
